\documentclass{article}
\usepackage{amsmath}
\begin{document}
\section*{The audit value of unresolved source identity}
There are $n$ distinct subjects, $g$ interchangeable complete readers,
$1\le g<n$, and one sparse reader that can verify exactly one of the
subjects. A pre-choice information cell leaves $k$ possible identities
for the sparse receipt. Choose subjects before resolving the identity
but assign them to readers afterward. Each reader verifies at most one
subject. Every choice in the distribution must admit an assignment in
every remaining role. The service value is the least inclusion
probability of any subject, with the controller's randomization private
and independent of the unrevealed role. A late assignment can use the
resolved role; an early decision cannot. This law is about available
qualified readers, not about how to obtain physical clearance or
enforce audit outcomes.

\paragraph{Exact value.}
Any set of at most $g$ subjects fits. A set of size $g+1$ fits if and
only if it contains the whole cell. If $g=n-1$, always choose all
subjects. Otherwise, if $k>g+1$, choose a uniform size-$g$ set and
obtain the sharp bound $g/n$. For $g<n-1$ and $k\le g+1$, write
$b=n-k$ and let
\[
q_*=\begin{cases}
g/(n-1), & b\ge g,\\
b(n-g)/(b+n(n-g-1)), & b<g .
\end{cases}
\]
Then the exact value is $p_*=(g+q_*)/n$. Here $0<q_*<1$; no
vanishing denominator is used because $k\le g+1<n$.

For any policy, extend smaller sets to size $g$ without lowering any
marginal. Excluding one cell member, every feasible choice contains
at most $g$ of the remaining $n-1$ subjects, so $p\le g/(n-1)$.
If the policy chooses a large set with probability $q$, total
expected inclusion is $g+q$. Its outsider count is at most
$b-q(n-g-1)$ in expectation. Thus $np\le g+q$ and
$bp\le b-q(n-g-1)$, yielding the displayed $p_*$ in the $b<g$
case by eliminating $q$. For attainment choose large sets with
probability $q_*$, always take all cell members and choose the other
$g+1-k$ uniformly outside. On small sets take an average of
\[
a=\frac{k(p_*-q_*)}{1-q_*}
\]
from the cell. Substitution gives $a=0$ if $b\ge g$ and $a=g-b$
otherwise: both are permitted small-set counts. Choose uniformly
within the two classes. Each cell member has inclusion $p_*$ and
the sum is $g+q_*$; outsider symmetry gives $p_*$ there too.
For $g=1$ this recovers $1/(n-1)$ exactly when the cell has at
most two possibilities. The original three-subject case needs just
one timely membership bit, not the exact role.

\paragraph{Queries before commitment.}
Suppose the only evidence action asks whether one named receipt is
the sparse one, returns a qualified yes or no immediately, and is
charged separately from the $g+1$ verification slots. A deterministic
policy that improves uniformly on the blind $g/n$ while $g<n-1$
must shrink every cell to at most $g+1$ candidates. On the all-no
path, each distinct query removes at most one, so $n-g-1$ queries
are necessary. Ask that many distinct questions, stopping on a yes;
this attains the bound. Exact identification needs $n-1$ in the
same model. Batch/group questions, indexes built for the current
manifest, absent or stale answers, and posterior-error guarantees
are not covered by this information lower bound. In particular no
transport-byte improvement follows when the full role and the cue
both fit in the same octet.

\paragraph{Source check.}
The sandbox replay uses separately backed SQLite archives from an
actual selected warehouse worker. It confirms the queried receipt
against both an exact manifest tuple and its keyed log effect in
one read-only transaction. Competing queries recover the full role
on the same archives. Query calls use independently initialized,
equally warmed connections and order permutation; VM progress counts
are local diagnostics for the stated SQLite build, not execution time
or a lower bound for all SQL plans. The single qualified point query
takes 24 callbacks in the positive case and 11 in either negative
case; the fixed-order adaptive full identification takes 24, 35,
or 22. The batch check takes 51. Role exclusivity and archive
availability are known in these cases. A stale or missing record is
not evidence that some other receipt exists unless the exactly-one
law has been established separately.

An independent matching-based linear program agrees on tested finite
cases; the proof above supplies the universal bound rather than the
floating-point check. This is a named reconstruction and quantitative
audit-policy result, not actual enforcement, measured runtime
advantage, an independent mission, or full R5/I1 acceptance.
\end{document}
